%-------------------------------------------------------%
%                         中文摘要
%-------------------------------------------------------%
\chapter*{摘\quad 要}
\cabstractpagestyle % 中文摘要的页眉页脚

在大数据、深度学习、人工智能与智能网联技术快速发展的背景下，针对城市公交系统普遍存在的串车、
车距不稳定等运行效率问题，本研究提出了一种基于深度强化学习（Deep Reinforcement Learning）的公交动态驻站控制模型。
针对传统驻站控制在实时响应、动态环境适应等方面的不足，创新性地构建了多智能体协同优化框架，
通过整合智能网联车辆数据与强化学习算法实现公交运行的实时动态驻站控制。

针对公交运营环境的动态性与复杂性，首先深入剖析公交的实际运行环境，将其中的核心要素进行抽象化处理，
构建相应模型。随后，对各类强化学习算法展开全面阐释与分析，经综合考量，采用演员-评论员（Actor - Critic）
算法框架来解决本研究面临的问题。选取多智能体近端策略优化算法（Multi-Agent Proximal Policy Optimization）
作为核心基础算法，在此基础上，进行算法优化改进，
% 通过分析公交的实际运行环境，对其中的核心要素进行抽象建模。同时，
% 对强化学习算法展开详尽阐述与分析，筛选出适配于本研究问题的算法框架Actor-Critic。
% 本研究选取多智能体近端策略优化算法（MAPPO）作为基础算法，
以此为依托构建了公交车队运行模型，设计了强化学习下公交运行模拟流程，并
通过分析影响公交系统车头时距稳定性的因素，提出采用公交驻站控制来提升公交运行稳定性。

本文将公交驻站控制建模为部分可观测的马尔可夫博弈（Markov Game）。通过
设计事件图（Event Graph）网络描述车辆到站的异步时序特征，
采用图注意力网络（Graph Attention Network）来学习事件图结构信息，
量化车辆间时空交互的影响强度，以此来补充马尔可夫博弈局部观测独立性约束造成无法考虑其他智能体
动作影响的弊端。基于此，提出改进的多智能体近端策略优化（Multi-Agent Proximal Policy Optimization）算法，
即在原始算法中引入图注意力网络，形成具有环境感知与协同决策能力的动态驻站控制模型。
在该模型中，将每辆公交车视为一个智能体，设计多维状态空间（包括反映车辆的运行稳定性的信息以及乘客需求信息）与
连续动作空间（驻站时长决策），并构建融合车头时距稳定性、乘客等待时间最小化的复合奖励函数。

本文设计了一条单向环形公交线路，将模型应用于该场景，通过选取对比基准模型
与评价指标对模型的实验结果进行分析，验证了模型的性能优势。为进一步验证模型的现实可行性，
将模型的应用场景扩展至常规公交线路，考虑公交运行中更多现实环境因素。
以成都市T34路公交为真实场景验证平台，
将模型的结果与传统的基于规则的驻站控制模型、数学优化控制模型以及基于原始的
MAPPO算法控制模型的结果进行对比，从微观和宏观层面多维度分析，
验证了本模型的有效性。数值结果显示，本文提出的基于改进MAPPO算法
的公交驻站控制模型相比于无控制模型（NH）、基于固定车头时距的控制模型（HH）、基于前向车头时距的控制模型（FH）、
基于后向车头时距的控制模型（BH）和基于原始MAPPO算法的控制模型，在乘客平均等待时间（AWT）指标上性能分别提升了 34.95\%、27.97\%、12.56\%、
20.00\%和7.84\%；在乘客平均旅行时间（ATT）指标上性能分别提升了
4.3\%、21.25\%、16.63\%、8.99\%和2.94\%；在车辆的平均占用离散度（AOD）指标上性能分别提升了 77.14\%、5.88\%、11.11\%、20.0\%
和11.11\%；在车辆的平均驻站时间（AHT）指标上，性能相比于基于固定车头时距的控制模型（HH）、
基于前向车头时距的控制模型（FH）、基于后向车头时距的控制模型（BH）和基于原始 MAPPO 算法的控制模型，
分别提升了 74.51\%、69.05\%、45.83\% 和27.78\%。
结果表明，本文提出的模型在保持公交系统车头时距稳定、减少公交串车、降低乘客出行时间成本、提高公交资源利用率以及提升
公交服务质量方面有良好表现。
% 实验表明，改进的MAPPO-GAT模型相
% 较于传统定时驻站控制方法，可将车头时距波动降低32.7%，乘客平均等待时间减少19.6%；与原始MAPPO算
% 法相比，优化效果提升11.3%，验证了模型在复杂动态环境下的鲁棒性。进一步以成都市T34路公交实际线路
% 为案例，通过整合GPS轨迹数据与站点客流信息，证实模型在真实路网中可有效缓解高峰时段的串
% 车现象，车距偏离度降低至4.2分钟以内，显著提升了公交服务可靠性。

% 本研究突破传统静态调度思维局限，首次将图神经网络与多智能体强化学习结合应用于公交动态控制领域，
% 为城市公交系统实时优化提供了可扩展的技术路径。未来研究可进一步探索多线路协同调度、混合交通场景
% 下的车路协同机制，以及基于车联网的分布式决策框架，推动公交系统智能化升级。
\vspace{1em}

\noindent \hangindent=3.9em \textbf{关键词：}强化学习，公交串车，智能体，驻站控制，异步特性